% GENERATED FILE. Edit canonical lessons, then run npm run generate:books.
% canonical-source-hash: fnv1a64:5c060b1251c99209
% canonical-lessons: TA-C532-kuttuccantai, TA-C532-malyuttam, TA-C532-cilampam, TA-C532-utai-kick, TA-C532-eri-throw

\chapter{Fighting Sports}
\label{ch:ta-fighting-sports}
\hlchaptermodality{532}

\begin{chapteropening}
\textbf{By the end of this chapter:} \emph{I can talk about boxing, wrestling and stick fencing.}

Complete the last lesson of chapter 532: I can talk about boxing, wrestling and stick fencing.
\end{chapteropening}

\section[\texorpdfstring{\ta{குத்துச்சண்டை} (kuttuccaṇṭai)}{குத்துச்சண்டை (kuttuccaṇṭai)}]{\ta{குத்துச்சண்டை} (kuttuccaṇṭai) — boxing}

\label{lesson:TA-C532-kuttuccantai}



\begin{quote}
Before the new one: say the Tamil for tennis, then the Tamil for kabaddi.
\end{quote}

\subsection*{You'll want to know: \ta{குத்துச்சண்டை}}
\textbf{\ta{குத்துச்சண்டை}} — \emph{kuttuccaṇṭai} — \textquotedblleft{}boxing\textquotedblright{}.

\begin{quote}
\ta{அவன்} \ta{குத்துச்சண்டையில்} \ta{முதல்} \ta{பரிசு} \ta{வென்றான்}.
\end{quote}

\begin{quote}
\emph{avaṉ kuttuccaṇṭaiyil mutal paricu veṉṟāṉ.}
\end{quote}

\begin{quote}
He won first prize in boxing.
\end{quote}

Word for word \textquotedblleft{}punch-fight\textquotedblright{}: \ta{குத்து} \textquotedblleft{}a punch\textquotedblright{} + \ta{சண்டை} \textquotedblleft{}a fight\textquotedblright{}.

\begin{cousinweb}
\textbf{\ta{குத்துச்சண்டை}} (\emph{kuttuccaṇṭai}) is a compound of the native \ta{குத்து} (\emph{kuttu}), \textquotedblleft{}to pierce, a thrust\textquotedblright{} (DEDR 1719), and \ta{சண்டை} (\emph{caṇṭai}), \textquotedblleft{}a fight\textquotedblright{}.
\end{cousinweb}

\begin{grammarlens}[title={What you've built}]
One more word for describing what you lived and felt.
\end{grammarlens}

\subsection*{Your turn}
\begin{itemize}
\raggedright
  \item \emph{Say it:} \emph{kuttuccaṇṭai}
  \item \emph{Say it:} \emph{kuttuccaṇṭai}, once more
  \item \emph{Say it:} \emph{avaṉ kuttuccaṇṭaiyil mutal paricu veṉṟāṉ.} — He won first prize in boxing.
  \item \emph{Recall:} say the Tamil for tennis, then the Tamil for kabaddi, then say \emph{kuttuccaṇṭai} again
\end{itemize}

\begin{tcolorbox}[breakable,colback=teal!4,colframe=teal!35!black,title={Before you move on}]
What does \ta{குத்துச்சண்டை} mean? (\textquotedblleft{}Boxing\textquotedblright{}.) Say it once more.
\end{tcolorbox}
\section[\texorpdfstring{\ta{மல்யுத்தம்} (malyuttam)}{மல்யுத்தம் (malyuttam)}]{\ta{மல்யுத்தம்} (malyuttam) — wrestling}

\label{lesson:TA-C532-malyuttam}



\begin{quote}
Before the new one: say the Tamil for kabaddi, then the Tamil for boxing.
\end{quote}

\subsection*{You'll want to know: \ta{மல்யுத்தம்}}
\textbf{\ta{மல்யுத்தம்}} — \emph{malyuttam} — \textquotedblleft{}wrestling\textquotedblright{}.

\begin{quote}
\ta{மல்யுத்தத்தில்} \ta{இருவரும்} \ta{சமமாகப்} \ta{போராடினார்கள்}.
\end{quote}

\begin{quote}
\emph{malyuttattil iruvarum camamākap pōrāṭiṉārkaḷ.}
\end{quote}

\begin{quote}
In the wrestling both fought evenly.
\end{quote}

A -\ta{அம்} noun: \ta{மல்யுத்தத்தில்}, in wrestling.

\begin{cousinweb}
\textbf{\ta{மல்யுத்தம்}} (\emph{malyuttam}) is borrowed from Sanskrit: \emph{malla}, \textquotedblleft{}a wrestler\textquotedblright{} (CDIAL 9907), and \emph{yuddha}, \textquotedblleft{}a battle\textquotedblright{} (CDIAL 10499). CDIAL itself marks \emph{malla} as taken from Dravidian.
\end{cousinweb}

\begin{grammarlens}[title={What you've built}]
One more word for describing what you lived and felt.
\end{grammarlens}

\subsection*{Your turn}
\begin{itemize}
\raggedright
  \item \emph{Say it:} \emph{malyuttam}
  \item \emph{Say it:} \emph{malyuttam}, once more
  \item \emph{Say it:} \emph{malyuttattil iruvarum camamākap pōrāṭiṉārkaḷ.} — In the wrestling both fought evenly.
  \item \emph{Recall:} say the Tamil for kabaddi, then the Tamil for boxing, then say \emph{malyuttam} again
\end{itemize}

\begin{tcolorbox}[breakable,colback=teal!4,colframe=teal!35!black,title={Before you move on}]
What does \ta{மல்யுத்தம்} mean? (\textquotedblleft{}Wrestling\textquotedblright{}.) Say it once more.
\end{tcolorbox}
\section[\texorpdfstring{\ta{சிலம்பம்} (cilampam)}{சிலம்பம் (cilampam)}]{\ta{சிலம்பம்} (cilampam) — stick fencing}

\label{lesson:TA-C532-cilampam}



\begin{quote}
Before the new one: say the Tamil for boxing, then the Tamil for wrestling.
\end{quote}

\subsection*{You'll want to know: \ta{சிலம்பம்}}
\textbf{\ta{சிலம்பம்}} — \emph{cilampam} — \textquotedblleft{}stick fencing\textquotedblright{}.

\begin{quote}
\ta{என்} \ta{மாமா} \ta{சிறுவர்களுக்குச்} \ta{சிலம்பம்} \ta{கற்றுக்கொடுக்கிறார்}.
\end{quote}

\begin{quote}
\emph{eṉ māmā ciṟuvarkaḷukkuc cilampam kaṟṟukkoṭukkiṟār.}
\end{quote}

\begin{quote}
My uncle teaches stick fencing to boys.
\end{quote}

A Tamil martial art fought with a long bamboo staff.

\begin{cousinweb}
\textbf{\ta{சிலம்பம்}} (\emph{cilampam}) is native Dravidian: DEDR 2564 gives \ta{சிலம்பம்} for the art of fencing with a quarterstaff, with a Malayalam cognate; where the name itself comes from is not settled.
\end{cousinweb}

\begin{grammarlens}[title={What you've built}]
One more word for describing what you lived and felt.
\end{grammarlens}

\subsection*{Your turn}
\begin{itemize}
\raggedright
  \item \emph{Say it:} \emph{cilampam}
  \item \emph{Say it:} \emph{cilampam}, once more
  \item \emph{Say it:} \emph{eṉ māmā ciṟuvarkaḷukkuc cilampam kaṟṟukkoṭukkiṟār.} — My uncle teaches stick fencing to boys.
  \item \emph{Recall:} say the Tamil for boxing, then the Tamil for wrestling, then say \emph{cilampam} again
\end{itemize}

\begin{tcolorbox}[breakable,colback=teal!4,colframe=teal!35!black,title={Before you move on}]
What does \ta{சிலம்பம்} mean? (\textquotedblleft{}Stick fencing\textquotedblright{}.) Say it once more.
\end{tcolorbox}
\section[\texorpdfstring{\ta{உதை} (utai)}{உதை (utai)}]{\ta{உதை} (utai) — to kick}

\label{lesson:TA-C532-utai-kick}



\begin{quote}
Before the new one: say the Tamil for wrestling, then the Tamil for stick fencing.
\end{quote}

\subsection*{You'll want to know: \ta{உதை}}
\textbf{\ta{உதை}} — \emph{utai} — \textquotedblleft{}to kick\textquotedblright{}.

\begin{quote}
\ta{அவன்} \ta{பந்தை} \ta{வேகமாக} \ta{உதைத்தான்}.
\end{quote}

\begin{quote}
\emph{avaṉ pantai vēkamāka utaittāṉ.}
\end{quote}

\begin{quote}
He kicked the ball hard.
\end{quote}

Past (I …): \ta{உதைத்தேன்}, \emph{utaittēṉ}.

What is kicked takes -\ta{ஐ}: \ta{பந்தை} \ta{உதை}. The noun is the same: \ta{ஒரு} \ta{உதை}, a kick.

\begin{cousinweb}
\textbf{\ta{உதை}} (\emph{utai}) is native Dravidian: DEDR 616 gives \ta{உதை} (\emph{utai}), \textquotedblleft{}to kick, to plant the foot against\textquotedblright{}.
\end{cousinweb}

\begin{grammarlens}[title={What you've built}]
One more word for describing what you lived and felt.
\end{grammarlens}

\subsection*{Your turn}
\begin{itemize}
\raggedright
  \item \emph{Say it:} \emph{utai}
  \item \emph{Say it:} \emph{utai}, once more
  \item \emph{Say it:} \emph{avaṉ pantai vēkamāka utaittāṉ.} — He kicked the ball hard.
  \item \emph{Recall:} say the Tamil for wrestling, then the Tamil for stick fencing, then say \emph{utai} again
\end{itemize}

\begin{tcolorbox}[breakable,colback=teal!4,colframe=teal!35!black,title={Before you move on}]
What does \ta{உதை} mean? (\textquotedblleft{}To kick\textquotedblright{}.) Say it once more.
\end{tcolorbox}
\section[\texorpdfstring{\ta{எறி} (eṟi)}{எறி (eṟi)}]{\ta{எறி} (eṟi) — to hurl; to throw}

\label{lesson:TA-C532-eri-throw}



\begin{quote}
Before the new one: say the Tamil for stick fencing, then the Tamil for to kick.
\end{quote}

\subsection*{You'll want to know: \ta{எறி}}
\textbf{\ta{எறி}} — \emph{eṟi} — \textquotedblleft{}to hurl; to throw\textquotedblright{}.

\begin{quote}
\ta{அவள்} \ta{கல்லை} \ta{ஏரியில்} \ta{எறிந்தாள்}.
\end{quote}

\begin{quote}
\emph{avaḷ kallai ēriyil eṟintāḷ.}
\end{quote}

\begin{quote}
She threw the stone into the lake.
\end{quote}

Past (I …): \ta{எறிந்தேன்}, \emph{eṟintēṉ}.

Close to \ta{வீசு} (\emph{vīcu}); \ta{எறி} stresses a hard throw at something.

\begin{cousinweb}
\textbf{\ta{எறி}} (\emph{eṟi}) is native Dravidian: DEDR 859 gives \ta{எறி} (\emph{eṟi}), \textquotedblleft{}to throw, to discharge, to fling\textquotedblright{}.
\end{cousinweb}

\begin{grammarlens}[title={What you've built}]
One more word for describing what you lived and felt.
\end{grammarlens}

\subsection*{Your turn}
\begin{itemize}
\raggedright
  \item \emph{Say it:} \emph{eṟi}
  \item \emph{Say it:} \emph{eṟi}, once more
  \item \emph{Say it:} \emph{avaḷ kallai ēriyil eṟintāḷ.} — She threw the stone into the lake.
  \item \emph{Recall:} say the Tamil for stick fencing, then the Tamil for to kick, then say \emph{eṟi} again
\end{itemize}

\begin{tcolorbox}[breakable,colback=teal!4,colframe=teal!35!black,title={Before you move on}]
What does \ta{எறி} mean? (\textquotedblleft{}To hurl; to throw\textquotedblright{}.) Say it once more.
\end{tcolorbox}
